\section{Singular injectivity and an Iitaka conversion}
\label{inj:injectivity}

The trivial-torus case of Theorem~\ref{inj:transfer} already converts
unbounded fixed-error sections into an anticanonical section. We give an
independent proof using singular injectivity and vanishing for all higher
direct images. On the Iitaka model $f:V\to Y$ constructed below, with
$V$ birational over $X$, the final interpolation retains
$f_*\mathcal O_V(K_{V/X})$ itself. The character-sensitive proof instead retains a corrected base
line with a fixed multiplier ideal. The resulting twisted-form conversion
will be applied after forms descend to a finite étale cover in
Section~\ref{har:subsec:formdescent}. The base metric need only have locally
bounded plurisubharmonic weights; the anticanonical metric upstairs
remains smooth and semipositive.

\begin{lemma}[Injectivity and base cohomology]\label{inj:leray}
Let $g:V\to Y$ be a morphism of smooth projective complex varieties, let $A$
be very ample on $Y$, and give $A$ a smooth metric with positive curvature
$\omega_A$.  Suppose a line bundle $N$ on $V$ has a singular Hermitian
metric $h$ such that
\[
 \mathcal I(h)=\mathcal O_V,
 \qquad i\Theta_h(N)\ge\epsilon g^*\omega_A
 \quad\text{for some }\epsilon>0.
\]
Then
\[
 H^i(Y,R^jg_*\mathcal O_V(K_V+N))=0
 \qquad(i>0,\ j\ge0).
\]
Disconnected smooth source varieties are allowed, with the hypotheses on
each component.
\end{lemma}

\begin{proof}
The input is the injectivity theorem of Fujino--Matsumura
\cite[Theorem A]{FM2016}.  On a compact K\"ahler manifold it applies to a
possibly singular coefficient metric and a smooth semipositive multiplying
line bundle $S$, provided the former curvature dominates a positive
multiple of the latter.  Multiplication by any nonzero section of $S$
is injective on canonical cohomology with the coefficient multiplier
ideal.  Here we use $S=g^*(mA)$; the comparison constant can be
$\epsilon/m$, and the ideal is trivial.

We include the reduction to base cohomology, following the slicing and
Leray argument in \cite[proof of Proposition 1.9, Step 3]{FM2016}.
Induct on $\dim Y$.  Choose $m$ large enough for Serre vanishing for all
the finitely many sheaves $R^jg_*\mathcal F$, where
$\mathcal F=\mathcal O_V(K_V+N)$.  A sufficiently general divisor
$C\in|mA|$ is smooth, has smooth inverse image $V_C$, contains no
associated point of these direct images, and contains the image of no
component of $V$.  Its restricted coefficient metric has trivial
multiplier ideal.  To see the last assertion without assuming analytic
singularities, pull the locally integrable metric coefficients to the
incidence variety of the basepoint-free system on $V$.  That incidence
variety is a projective-space bundle over $V$.  Fubini's theorem gives
integrability on almost every smooth member, simultaneously on a finite
coordinate cover.  This full-measure condition meets the required
algebraic open conditions.  Restricting the curvature inequality on
these members gives the same hypothesis over $C$.

Adjunction and multiplication by the cutting section give
\[
 0\longrightarrow\mathcal F\longrightarrow
 \mathcal F\otimes g^*\mathcal O_Y(mA)\longrightarrow
 \mathcal O_{V_C}(K_{V_C}+N|_{V_C})\longrightarrow0.
\]
The corresponding long direct-image sequence breaks into short exact
sequences: multiplication by the cutting equation is injective on every
$R^jg_*\mathcal F$, by the associated-point choice.  The induction
hypothesis on $C$ and Serre vanishing for the middle terms therefore
give $H^i(Y,R^jg_*\mathcal F)=0$ for $i\ge2$.

The Leray spectral sequence now has only columns zero and one.  For
each $j$ it identifies $H^1(Y,R^jg_*\mathcal F)$ with the first
filtration subspace of $H^{j+1}(V,\mathcal F)$.  After the twist by
$g^*(mA)$ this subspace maps to zero, since all positive base cohomology
of the twisted direct images vanishes.  Functoriality of the Leray
filtration and Fujino--Matsumura injectivity imply that the original
subspace is zero.  This proves the remaining degree $i=1$.
\end{proof}

\begin{proposition}[Conversion on an Iitaka base]\label{inj:ordinary-conversion}
Let $X$ be smooth connected projective over $\mathbb C$ and let
$L=-K_X$ carry a smooth semipositive Hermitian metric.  Suppose that
$F_0$ is a pseudoeffective line bundle and, for an unbounded set
$S\subset\mathbb Z_{>0}$, there are effective integral divisors
\[
 N_s\sim sL-F_0\qquad(s\in S).
\]
Then $H^0(X,mL)\ne0$ for some integer $m>0$.
\end{proposition}

\begin{proof}
Fix $a<b$ in $S$.  The affine relations
\begin{equation}\label{inj:affine-divisors}
 (s-a)N_b\sim(s-b)N_a+(b-a)N_s\qquad(s>b)
\end{equation}
will make the fiberwise section direction unique.  Put
$e=\kappa(X,N_b)$.  If $e=0$, the two effective divisors in
\eqref{inj:affine-divisors} are equal.  Letting $s$ tend to infinity
gives $N_b\ge N_a$, so $N_b-N_a\sim(b-a)L$ is effective.  If
$e=\dim X$, $bL=N_b+F_0$ is big, which also proves the assertion.
We treat $0<e<\dim X$.

Resolve a multiple system of $N_b$ attaining its Iitaka dimension and
take its Stein factorization.  After modifying the base, flattening
the main component, normalizing, and resolving, we have
\[
 V\longrightarrow W\longrightarrow X,
 \qquad f:V\to Y,
 \qquad f_W:W\to Y.
\]
Here $Y,V$ are smooth projective, $W$ is normal projective and birational
to $X$, $f_W$ is equidimensional, and $f$ has connected fibers.  The
normalization is finite and preserves the fiber-dimension bound from
flattening; the resolution $V\to W$ need not be equidimensional over
$Y$.  Equivalently the flat main component is obtained from the closure
of the family of generic fibers in the projective Hilbert scheme
\cite[Theorem~3.2]{Grothendieck1961Hilbert}.  Write
$\pi:V\to X$, $L'=\pi^*L$, and $E=K_{V/X}\ge0$.
For a very ample $A$ on $Y$ and some integer $t>0$,
\[
 t\pi^*N_b-f^*A\sim G\ge0.
\]
The initial moving polarization stays big after base modification, so
it dominates the chosen $A$ after increasing $t$.

For a geometric generic fiber $F$,
\begin{equation}\label{inj:rigidity}
 \kappa\bigl(F,(\pi^*N_b+E)|_F\bigr)=0.
\end{equation}
Indeed adding an effective exceptional divisor changes no global
sections of multiples of a pullback from the smooth $X$.  If two
sections on the generic fiber were independent, a sufficiently large
ample base twist would globalize two such sections.  Absorb that twist
using $t\pi^*N_b\ge f^*A$, up to linear equivalence.  Their nonconstant
fiberwise ratio, together with the ratios giving the base, would yield
image dimension at least $e+1$ for a multiple of $\pi^*N_b+E$.
The geometric generic fiber is integral because the Stein field is
relatively algebraically closed in characteristic zero.  Generic base
change justifies the globalization over the original function field.
This contradiction proves \eqref{inj:rigidity}.  The same assertion
holds without $E$.

Restriction of \eqref{inj:affine-divisors} to $F$ is therefore equality
of effective divisors.  Unboundedness of $S$ shows
$(\pi^*(N_b-N_a))|_F\ge0$.  If $d=b-a$, a rational section $z$ of
$dL$ with divisor $N_b-N_a$ is consequently regular on $F$.  For every
$k\ge0$,
\begin{equation}\label{inj:rank-one-all-k}
 h^0\bigl(F,(E+kdL')|_F\bigr)=1.
\end{equation}
The indicated effective section gives the lower bound.  For $k\ge1$,
multiply by the section of $k\pi^*N_a+(k-1)E$ to inject into the
sections of $k(\pi^*N_b+E)|_F$.  For $k=0$, multiply by the section
of $\pi^*N_b|_F$.  Equation \eqref{inj:rigidity} gives both upper
bounds.  This explicitly includes the zero exponent.

On the normal equidimensional model $W$, take for every prime $P\subset Y$
the minimum
\[
 c_P=\min_{Q\mapsto P}
       \frac{\operatorname{ord}_Q(z)}{
       \operatorname{ord}_Q(f_W^*P)}.
\]
Every vertical prime maps to a base prime and every denominator is
positive.  There are only finitely many nonzero $c_P$.  Subtracting
$f_W^*(\sum c_PP)$ leaves an effective rational Cartier divisor,
because the horizontal divisor of $z$ was already effective.
Clear denominators, replacing $d,z$ by a common power.  We obtain
\begin{equation}\label{inj:base-decomposition}
 dL'\sim H+f^*B,
 \qquad H\ge0,
\end{equation}
with $H$ pulled back from an effective Cartier divisor on $W$ and
$B$ integral on $Y$.  Above each base prime at least one component
has coefficient zero in $H$.

Proposition~\ref{har:prop:maximum} now applies with the ordinary line $P=L$, the rational section $z$, and the normal equidimensional model $W$. Equation~\eqref{inj:rank-one-all-k} supplies rank one for every $k\ge0$; the vertical minima giving \eqref{inj:base-decomposition} agree with the proposition's normalization. Hence $B$ has a continuous semipositive metric. In a local frame $b$, with $\sigma=s_Hf^*b$ and $h_0=\pi^*h_L$, its weight is
\begin{equation}\label{inj:max-weight}
 \phi_B=-\log\max_{V_y}|\sigma|_{h_0^d}^2.
\end{equation}
Only local boundedness of this metric is needed in the following injectivity argument.

It remains to use the zero exponent rather than merely asymptotic
positivity.  Since $F_0$ is pseudoeffective, the ample domination above
gives $nL'-f^*A$ pseudoeffective for some $n>0$.  Its positive current
yields a singular metric on $L'$ with curvature dominating a positive
multiple of $f^*\omega_A$.  Mix its weights with the smooth semipositive
weights using a sufficiently small positive coefficient.  Skoda's
small-exponent integrability \cite[Lemma~5.6(a)]{Demailly2011}, in
the squared-norm convention $e^{-\phi}$, and a finite covering of $V$ make the
new multiplier ideal trivial while retaining a strictly positive base
curvature coefficient.  No analytic-singularity assumption is used.
For $k\ge0$ put $C_k=L'+kf^*B$ and add $kf^*\phi_B$ to the coefficient
weights. Local boundedness keeps the multiplier ideal trivial and
preserves the positive base-curvature lower bound. Since
$K_V+C_k=E+kf^*B$, Lemma~\ref{inj:leray} and the projection formula give
\[
 H^i\bigl(Y,f_*\mathcal O_V(E)\otimes\mathcal O_Y(kB)\bigr)=0
 \quad(i>0,\ k\ge0).
\]
Riemann--Roch makes the Euler characteristic a polynomial in $k$.
At $k=0$ it equals $h^0(V,E)=1$.  It is therefore nonzero at some
positive integer $k$.  A section of $E+kf^*B$, multiplied by $s_H^k$,
gives a section of $E+kdL'$.  Since
$\pi_*\mathcal O_V(E)=\mathcal O_X$, it yields a nonzero section of
$kdL$ on $X$.
\end{proof}

\begin{corollary}[Twisted-form input]\label{inj:forms-input}
Under the smooth projective and smooth anticanonical semipositivity
hypotheses of Proposition~\ref{inj:ordinary-conversion}, suppose
$H^0(X,\Omega_X^p\otimes mL)\ne0$ for an unbounded set of positive
integers $m$ and a fixed $p\ge0$.  Then a positive power of $L$ has
a nonzero section.
\end{corollary}
\begin{proof}
For $p=0$ the assertion is immediate.  Otherwise take the stable generic
span of the directions of the twisted forms in $\Omega_X^p$, of rank
$r>0$, and saturate its determinant in
$\bigwedge^r\Omega_X^p\subset(\Omega_X^1)^{\otimes pr}$.
A saturated rank-one subsheaf of a locally free sheaf on smooth $X$
is reflexive and hence invertible; call it $M$.  Wedge an arbitrarily
far-tail member with suitable $r-1$ members of a fixed generic basis.
This gives nonzero sections of $M+sL$ for unbounded sums $s$.
The cotangent-tensor generic-nefness theorem
\cite[Theorem 4.1]{LMPTX2023} applies because $-K_X$ is nef and
proves that $-M$ is pseudoeffective.  Apply the proposition with
$F_0=-M$ and the divisors of these sections.
\end{proof}

For the direct twisted-form descent in
Section~\ref{har:subsec:formdescent}, the forms in this
corollary must first descend to a smooth projective finite \'etale
cover of the original variety.  The conversion asserts ordinary
sections on that cover.  A norm then descends a section of a further
positive power.  It asserts neither equivariance of the conversion nor
a uniform bound on the positive exponent.
